% Generated draft: requires mathematical and source review.
\documentclass[11pt]{article}
\usepackage[margin=25mm]{geometry}
\usepackage{fontspec}
\setmainfont{Latin Modern Roman}
\usepackage{amsmath,amssymb,mathtools,unicode-math}
\setmathfont{Latin Modern Math}
\usepackage{enumitem,needspace,xurl,hyperref}
\setlength{\parindent}{0pt}
\setlength{\parskip}{5pt}
\setlength{\emergencystretch}{3em}
\newcommand{\N}{\mathbb{N}}
\newcommand{\Z}{\mathbb{Z}}
\newcommand{\Q}{\mathbb{Q}}
\newcommand{\R}{\mathbb{R}}
\newcommand{\C}{\mathbb{C}}
\newcommand{\F}{\mathbb{F}}
\newcommand{\A}{\mathbb{A}}
\newcommand{\PP}{\mathbb{P}}
\newcommand{\E}{\mathbb{E}}
\newcommand{\eps}{\varepsilon}
\newcommand{\dd}{\,\mathrm{d}}
\begin{document}
\section*{UnsolvedMath ID: 30003968}
\textbf{Source problem code:} catalogue entry
\subsection{Definitions and mathematical statement}
Let \(n_1,n_2\) denote the two population densities in the supplied tumor-invasion model, on \(\R^d\) for \(t\geq0\). Set
\[
n=n_1+n_2,\qquad p=n^\gamma,\qquad \gamma>1,
\]
and let \(F_1,F_2,G_1,G_2\) be pressure-dependent growth terms, with
\[
F=F_1+F_2,\qquad G=G_1+G_2.
\]
The system to be satisfied, in a weak sense that the record does not specify, is
\[
\begin{aligned}
\partial_t n_1-\operatorname{div}(n_1\nabla p)
 &=n_1F_1(p)+n_2G_1(p),\\
\partial_t n_2-\operatorname{div}(n_2\nabla p)
 &=n_1F_2(p)+n_2G_2(p)
\end{aligned}
\qquad\text{on }\R^d\times(0,\infty).
\]
The supplied homeostatic pressure assumption is
\[
F(p)\leq0,\qquad G(p)\leq0\qquad\text{for every }p\geq P_H.
\]
\textbf{Question as formulated in the supplied statement.} Can an existence theory for weak solutions of this system be established while retaining the displayed homeostatic pressure assumption and dispensing with the growth compatibility condition
\[
\sup_{0\leq p\leq P_H}
\frac{|F(p)-G(p)|^2}{p^{1/\gamma}}\leq C_H,
\]
that is, without imposing such a finite uniform bound?

\textbf{Unresolved qualifications.} This is an incomplete existence question, rather than a fully quantified theorem. The record supplies no initial data or admissible class of initial data, no precise weak-solution definition, no regularity or sign assumptions on the individual growth terms, and no integrability or behavior at spatial infinity. The permitted dimensions, the range of \(P_H\), and the permitted dependence of \(C_H\) on parameters are not specified. Although \(n_1,n_2\) are called densities, their nonnegativity is not explicitly included among the supplied hypotheses; the admissible range of \(n\) must also make \(n^\gamma\) meaningful. These missing specifications prevent a precise assertion of existence for every admissible datum.

The displayed quotient is undefined at \(p=0\) as written. A supremum over \(0<p\leq P_H\), or an explicitly prescribed extension at zero, would require a clarification not supplied by the record. Furthermore, the supplied literature assessment describes removal of the homeostatic pressure assumption, whereas the mathematical statement asks to remove the growth compatibility bound while retaining the homeostatic inequalities. These targets must be distinguished pending source review.
\subsection{Short English statement}
The supplied question asks whether weak solutions of a two-population tumor-invasion system with pressure \(p=(n_1+n_2)^\gamma\), \(\gamma>1\), can be constructed without the stated bound on the difference of the combined growth terms \(F\) and \(G\). The displayed formulation retains their nonpositivity above the homeostatic pressure. Missing hypotheses, the undefined quotient at zero, and a conflicting description in the supplied literature assessment require human review.
\subsection{Sources}
\begin{itemize}
\item Local UnsolvedMath registry, \texttt{data/unsolvedmath/problems.json}, supplied record with permanent UnsolvedMath ID \(30003968\), source problem code \texttt{catalogue entry}, title \emph{Weak Tumor-Invasion Solutions Without Growth Compatibility}. The supplied collection is \emph{Oberwolfach Reports Open Problems}; the proposed year is 2018, and no individual proposer is identified.
\item Supplied source citation: \emph{Differential Equations arising from Organising Principles in Biology} (2019), Oberwolfach Reports, pp.\ 2651--2701, DOI \texttt{10.4171/owr/2018/43}. \url{https://doi.org/10.4171/owr/2018/43}. The record describes the reports as documenting mathematical workshops and their open-problem discussions. Its verification note says that the official report was checked and that the reconstruction restores the population equations, pressure law, growth terms, homeostatic assumption, and compatibility condition. This is supplied verification metadata, not an independent verification performed for this draft.
\item Additional supplied literature source, labelled \emph{Two species tissue growth}: \url{https://arxiv.org/abs/1809.01867}. No authors or further bibliographic details are supplied in the record.
\item Supplied literature assessment, checked on \texttt{2026-08-21}: the record labels the problem open, describes the model record as an incomplete fragment, states that the source calls removal of the homeostatic pressure assumption open, and says that later work proves existence only under assumptions. This is supplied dated curation derived from the OWR material, not an independently verified literature assessment. The record itself requests independent verification of time-sensitive status and cited claims.
\end{itemize}
\end{document}
